% ============================================================
%  Extended Abstract — LEAF 2026
%  Format: A4, 25 mm top/left/right, 30 mm bottom
%          Times New Roman 11 pt, single spacing, justified
% ============================================================
\documentclass[11pt, a4paper]{article}

% ── Geometry ─────────────────────────────────────────────────
\usepackage[a4paper,
            top    = 25mm,
            left   = 25mm,
            right  = 25mm,
            bottom = 30mm]{geometry}

% ── Fonts & Encoding ─────────────────────────────────────────
\usepackage[T1]{fontenc}
\usepackage[utf8]{inputenc}
\usepackage{times}           % Times New Roman

% ── Math ─────────────────────────────────────────────────────
\usepackage{amsmath, amssymb}

% ── Graphics & Tables ────────────────────────────────────────
\usepackage{graphicx}
\usepackage{booktabs}
\usepackage{caption}
\captionsetup{font=small, labelfont=bf, justification=centering}

% ── Spacing & Layout ─────────────────────────────────────────
\usepackage{setspace}
\singlespacing
\setlength{\parindent}{0pt}
\setlength{\parskip}{6pt}

% ── Hyperlinks ───────────────────────────────────────────────
\usepackage[hidelinks]{hyperref}

% ── Enumeration ──────────────────────────────────────────────
\usepackage{enumitem}
\setlist[enumerate]{leftmargin=*, labelwidth=1.8em, align=left,
                    itemsep=2pt, parsep=0pt, topsep=2pt}

% ── Section: no numbering, ABSTRACT style ────────────────────
\usepackage{titlesec}
\titleformat{\section}[block]
  {\normalfont\normalsize\bfseries\centering}
  {}{0pt}{}
\titlespacing*{\section}{0pt}{12pt}{6pt}

% ============================================================
\begin{document}
\pagestyle{empty}

% ── Title ────────────────────────────────────────────────────
\begin{center}
  {\large\bfseries
    STATIC and Dynamic Biplot Approaches for Socio-Economic Spatio-Temporal Analysis \\
    (Case Studies: Province of Banten, DKI Jakarta, and West Java 2020--2024)
  \par}

  \vspace{10pt}

  % ── Authors ────────────────────────────────────────────────
  {\normalsize\bfseries
    Ammar Hanafi\textsuperscript{1,*}
    and
    Yekti Widyaningsih\textsuperscript{1}
  \par}

  \vspace{6pt}

  % ── Affiliations ───────────────────────────────────────────
  {\small\itshape
    \textsuperscript{1}Department of Mathematics, Faculty of Mathematics and
    Natural Sciences,\\
    Universitas Indonesia, Kampus UI Depok, Depok 16424, Indonesia
  \par}

  \vspace{4pt}

  {\small
    *Corresponding author:
    \href{mailto:ammar.hanafi@ui.ac.id}{ammar.hanafi@ui.ac.id}
  \par}
\end{center}

\vspace{8pt}

% ── Abstract ─────────────────────────────────────────────────
\section*{ABSTRACT}

This study applies the \textit{STATIS} (\textit{Structuration des Tableaux
\`{a} Trois Indices de la Statistique}) method to jointly analyse 11
socio-economic indicators across 41 regencies and cities in Banten, DKI
Jakarta, and West Java over 2020--2024, capturing the structural impact of
the COVID-19 pandemic and an uneven post-pandemic recovery.
Demographic weights derived from population figures ensure a metrically
consistent comparison space.
Inter-structure evaluation via $RV$ coefficients confirms exceptionally high
longitudinal consistency (all pairwise $RV > 0.90$; $p < 0.05$),
demonstrating that the regional development hierarchy remained structurally
stable despite pandemic fluctuations.
The STATIS compromise explains $71.88\%$ of total variance: Dimension~1
captures socio-economic capacity (correlated with GRDP and HDI), while
Dimension~2 reflects labour-market dynamics (open unemployment rate).
Dynamic Biplot trajectories identify three regional patterns: V-shaped
recoveries in service-based metropolitan centres, stagnant or declining
trajectories in southern regencies, and progressive growth in buffer-zone
cities.
K-Means clustering ($k=3$) partitions the 41 regions into
\textit{Advanced} (13), \textit{Developing} (16), and
\textit{Underdeveloped} (12) clusters, with strong alignment to UNDP HDI
normative categories.
These results indicate that regional inequality is structurally persistent,
necessitating targeted rather than uniform policy intervention.

\vspace{8pt}

% ── Figures (side by side) ────────────────────────────────────
\begin{minipage}[t]{0.48\linewidth}
  \centering
  \IfFileExists{D:/Skripsi/my/output_real/scenario1/compromise_map.png}{%
    \includegraphics[width=\linewidth]{D:/Skripsi/my/output_real/scenario1/compromise_map.png}%
  }{%
    \fbox{\parbox{\linewidth}{\centering\vspace{2.8cm}
      \small[compromise\_map.png]\vspace{2.8cm}}}%
  }
  \captionof{figure}{Spatial distribution of K-Means clusters ($k=3$)
  across 41 regencies/cities ($k=3$, 2022).}
\end{minipage}%
\hfill%
\begin{minipage}[t]{0.48\linewidth}
  \centering
  \IfFileExists{D:/Skripsi/my/output_real/scenario1/compromise_scatter.png}{%
    \includegraphics[width=\linewidth]{D:/Skripsi/my/output_real/scenario1/compromise_scatter.png}%
  }{%
    \fbox{\parbox{\linewidth}{\centering\vspace{2.8cm}
      \small[compromise\_scatter.png]\vspace{2.8cm}}}%
  }
  \captionof{figure}{Compromise biplot scores with cluster overlays
  and trajectory arrows ($k=3$).}
\end{minipage}

\vspace{10pt}

% ── Cluster profile table ─────────────────────────────────────
\section*{CLUSTER PROFILES (STATIS Compromise, $k=3$, 2022)}

\noindent
Table~\ref{tab:cluster} summarises the mean indicator profiles per cluster
under the STATIS Compromise (ACT-Classic, $k=3$, 2022).
GRDP and mean years of schooling increase monotonically from Cluster~1 to
Cluster~3, while poverty rates follow the inverse pattern.
Unemployment (TPT) is notably higher in Clusters~2 and~3 than in
Cluster~1, consistent with a structural labour-market paradox in
service-led metropolitan economies.

\vspace{6pt}

\begin{center}
{\small
\begin{tabular}{lccc}
\toprule
\textbf{Indicator} & \textbf{Cluster 1} & \textbf{Cluster 2} & \textbf{Cluster 3} \\
                   & \textit{Underdeveloped} & \textit{Developing} & \textit{Advanced} \\
                   & ($n=12$) & ($n=16$) & ($n=13$) \\
\midrule
GRDP (billion IDR) & $27{,}370 \pm 18{,}004$ & $58{,}076 \pm 50{,}787$ & $215{,}387 \pm 158{,}352$ \\
Poverty rate (\%)  & $9.79 \pm 2.45$  & $8.85 \pm 2.44$  & $4.68 \pm 1.39$  \\
Mean yrs schooling & $7.70 \pm 0.69$  & $8.64 \pm 1.02$  & $11.14 \pm 0.59$ \\
Unemployment (\%)  & $6.55 \pm 2.60$  & $8.44 \pm 1.06$  & $8.22 \pm 1.69$  \\
\bottomrule
\end{tabular}
\captionof{table}{Mean $\pm$ SD of key indicators per cluster (Kompromi
STATIS, Scenario~1, $k=3$, reference year 2022).}
\label{tab:cluster}
}
\end{center}

\vspace{8pt}

% ── Keywords ─────────────────────────────────────────────────
\noindent
\textbf{Keywords}\quad:\quad
Interactive Visualization; K-Means Clustering; Longitudinal Data; PCA; Three-Way Data Analysis

\vspace{10pt}

% ── References ───────────────────────────────────────────────
\noindent\textbf{References}
\begin{enumerate}[label=\arabic*.]

\item C. Lavit, Y. Escoufier, R. Sabatier, and P. Traissac,
  ``The ACT (STATIS method),''
  \textit{Computational Statistics \& Data Analysis},
  vol.~18, no.~1, pp.~97--119, 1994.

\item H. Abdi, L. J. Williams, D. Valentin, and M. Bennani-Dosse,
  ``STATIS and DISTATIS: Optimum multitable principal component
  analysis and three way metric multidimensional scaling,''
  \textit{WIREs Computational Statistics}, vol.~4, no.~2,
  pp.~124--167, 2012.

\item Badan Pusat Statistik,
  \textit{Indeks Pembangunan Manusia 2020--2024},
  BPS-Statistics Indonesia, Jakarta, 2024.

\item J. MacQueen, ``Some methods for classification and analysis of
  multivariate observations,'' in \textit{Proc.\ 5th Berkeley Symp.\
  Mathematical Statistics and Probability}, vol.~1,
  University of California Press, 1967, pp.~281--297.

\end{enumerate}

\end{document}
